% --- Archon physics formalization source begin ---
% archon:physics
% archon:covers PhyXMiniProblems/problem_phyx_mini_0636.lean
% archon:source-report reports/phyx_mini/problem_phyx_mini_0636.source.json

\chapter{Physics problem phyx\_mini\_0636}
\label{ch:PhyXMiniProblems_problem_phyx_mini_0636}

\paragraph{Problem source.}
\#\# Physical scenario

The cesium isotope \$\textasciicircum{}\{137\}\$Cs, with \$Z=55\$, is radioactive and decays by beta decay. A beta particle is observed in the laboratory with a kinetic energy of \$300\textbackslash{};keV\$.

\#\# Question

With what kinetic energy was the beta particle ejected from the \$12.4\textbackslash{};fm\$-diameter nucleus?

\#\# Answer choices

- A: A: \$12.7\textbackslash{};MeV\$

- B: B: \$13.1\textbackslash{};MeV\$

- C: C: \$13.4\textbackslash{};MeV\$

- D: D: \$13.8\textbackslash{};MeV\$

\#\# Figure caption (auxiliary; use the image as primary evidence)

The image features a diagram of a physical interaction involving a cesium nucleus and a negatively charged particle.

**Left Side:**
- **Text:**
  - "Before: \textbackslash{}(K\_i\textbackslash{}), \textbackslash{}(U\_i\textbackslash{})"
  - "Cesium nucleus"
  - "R = 6.2 fm"
  - "q\_1 = 55e"
  - "q\_2 = -e"
- **Diagram:**
  - A circle representing the cesium nucleus with "+" signs inside to indicate positive charge.
  - A smaller circle with a "-" sign representing a negatively charged particle near the nucleus.
  - An arrow pointing away from the nucleus shows direction of motion.

**Right Side:**
- **Text:**
  - "After: \textbackslash{}(K\_f = 300 \textbackslash{}, \textbackslash{}text\{keV\}\textbackslash{})"
  - "U\_f = 0 \textbackslash{}, \textbackslash{}text\{keV\}"
- **Diagram:**
  - A smaller circle with a "-" sign and an arrow pointing to the right, indicating the particle's movement away from the nucleus.

The image illustrates the initial and final states of a charged particle interacting with a cesium nucleus, highlighting changes in kinetic and potential energy.

\#\# Dataset metadata

Nuclear Physics; Multi-Formula Reasoning, Physical Model Grounding Reasoning

\paragraph{Recorded answer/context.}
B: B: \$13.1\textbackslash{};MeV\$

\paragraph{Figure/image path.}
/root/proposal\_for\_physic/science-mango/phyx\_mini\_run/phyx\_data/test\_image/636.png

\paragraph{Formalization target.}
create a compiling Lean file with sorry bodies at `PhyXMiniProblems/problem\_phyx\_mini\_0636.lean`. The Lean declarations must preserve the physical quantities, dimensions or dimensional roles, figure labels, governing-law hypotheses, and final relation expressed by this problem.
Use Mathlib/Physlib names found through LeanExplore where available. If a physics API is missing, introduce faithful local abstractions rather than scalar placeholder aliases.

\begin{theorem}[Physics formalization target]
\label{thm:physics:phyx_mini_0636:target}
\lean{PhyXMiniProblems.ProblemPhyXMini0636.problem_phyx_mini_0636}
\uses{def:physics:phyx-mini-0636:phyxminiproblems-problemphyxmini0636-lengthinmeters, def:physics:phyx-mini-0636:phyxminiproblems-problemphyxmini0636-coulombconstantinjoulemeterspercoulombsquared, def:physics:phyx-mini-0636:phyxminiproblems-problemphyxmini0636-elementarychargeincoulombs, def:physics:phyx-mini-0636:phyxminiproblems-problemphyxmini0636-megaelectronvoltinjoules, def:physics:phyx-mini-0636:phyxminiproblems-problemphyxmini0636-energyinmegaelectronvolts, def:physics:phyx-mini-0636:phyxminiproblems-problemphyxmini0636-cesiumbetaejectionsetup, def:physics:phyx-mini-0636:phyxminiproblems-problemphyxmini0636-matchescesium137betadecayscenario, def:physics:phyx-mini-0636:phyxminiproblems-problemphyxmini0636-matchessuppliedcesiumbetafigure, def:physics:phyx-mini-0636:phyxminiproblems-problemphyxmini0636-usessuppliedfigurecalibration, def:physics:phyx-mini-0636:phyxminiproblems-problemphyxmini0636-hasnuclearsurfacegeometry, def:physics:phyx-mini-0636:phyxminiproblems-problemphyxmini0636-hasphysicalcesiumbetaparameters, def:physics:phyx-mini-0636:phyxminiproblems-problemphyxmini0636-usesvacuumcoulombconstant, def:physics:phyx-mini-0636:phyxminiproblems-problemphyxmini0636-satisfiesidealelectrostaticejectionlaws, def:physics:phyx-mini-0636:phyxminiproblems-problemphyxmini0636-matchesdisplayedejectionenergy, def:physics:phyx-mini-0636:phyxminiproblems-problemphyxmini0636-isuniquematchingdisplayedejectionenergy}
The assigned autoformalize agent should translate this physics problem into Lean declarations in the covered file, with theorem and lemma proof bodies written as `by sorry`.
\end{theorem}
\begin{proof}
This is an autoformalization task, not a proof task. Produce statements that can later be proved by the physics prover without weakening the physical model.
\end{proof}

% --- Archon declaration topology begin ---
\section{Declaration topology}
\begin{definition}[coulomb Constant Dimension]
  \label{def:physics:phyx-mini-0636:phyxminiproblems-problemphyxmini0636-coulombconstantdimension}
  \lean{PhyXMiniProblems.ProblemPhyXMini0636.coulombConstantDimension}
  The dimension of Coulomb's constant, equivalently J m / C²
\end{definition}
\begin{proof}
  This is the declared model definition of coulomb Constant Dimension.
\end{proof}
\begin{definition}[Length Quantity]
  \label{def:physics:phyx-mini-0636:phyxminiproblems-problemphyxmini0636-lengthquantity}
  \lean{PhyXMiniProblems.ProblemPhyXMini0636.LengthQuantity}
  A nonnegative, unit-independent physical length
\end{definition}
\begin{proof}
  This is the declared model definition of Length Quantity.
\end{proof}
\begin{definition}[Signed Charge Quantity]
  \label{def:physics:phyx-mini-0636:phyxminiproblems-problemphyxmini0636-signedchargequantity}
  \lean{PhyXMiniProblems.ProblemPhyXMini0636.SignedChargeQuantity}
  A signed, unit-independent electric charge
\end{definition}
\begin{proof}
  This is the declared model definition of Signed Charge Quantity.
\end{proof}
\begin{definition}[Coulomb Constant Quantity]
  \label{def:physics:phyx-mini-0636:phyxminiproblems-problemphyxmini0636-coulombconstantquantity}
  \lean{PhyXMiniProblems.ProblemPhyXMini0636.CoulombConstantQuantity}
  \uses{def:physics:phyx-mini-0636:phyxminiproblems-problemphyxmini0636-coulombconstantdimension}
  A nonnegative, dimensionful value of Coulomb's constant
\end{definition}
\begin{proof}
  This is the declared model definition of Coulomb Constant Quantity.
\end{proof}
\begin{definition}[length Readout]
  \label{def:physics:phyx-mini-0636:phyxminiproblems-problemphyxmini0636-lengthreadout}
  \lean{PhyXMiniProblems.ProblemPhyXMini0636.lengthReadout}
  \uses{def:physics:phyx-mini-0636:phyxminiproblems-problemphyxmini0636-lengthquantity}
  Read a physical length as a real number in a selected length unit
\end{definition}
\begin{proof}
  This is the declared model definition of length Readout.
\end{proof}
\begin{definition}[length In Meters]
  \label{def:physics:phyx-mini-0636:phyxminiproblems-problemphyxmini0636-lengthinmeters}
  \lean{PhyXMiniProblems.ProblemPhyXMini0636.lengthInMeters}
  \uses{def:physics:phyx-mini-0636:phyxminiproblems-problemphyxmini0636-lengthquantity, def:physics:phyx-mini-0636:phyxminiproblems-problemphyxmini0636-lengthreadout}
  Read a physical length in coherent-SI metres
\end{definition}
\begin{proof}
  This is the declared model definition of length In Meters.
\end{proof}
\begin{definition}[length In Femtometers]
  \label{def:physics:phyx-mini-0636:phyxminiproblems-problemphyxmini0636-lengthinfemtometers}
  \lean{PhyXMiniProblems.ProblemPhyXMini0636.lengthInFemtometers}
  \uses{def:physics:phyx-mini-0636:phyxminiproblems-problemphyxmini0636-lengthquantity, def:physics:phyx-mini-0636:phyxminiproblems-problemphyxmini0636-lengthreadout}
  Read a physical length in femtometres
\end{definition}
\begin{proof}
  This is the declared model definition of length In Femtometers.
\end{proof}
\begin{definition}[signed Charge In Coulombs]
  \label{def:physics:phyx-mini-0636:phyxminiproblems-problemphyxmini0636-signedchargeincoulombs}
  \lean{PhyXMiniProblems.ProblemPhyXMini0636.signedChargeInCoulombs}
  \uses{def:physics:phyx-mini-0636:phyxminiproblems-problemphyxmini0636-signedchargequantity}
  Read a signed physical charge in coherent-SI coulombs
\end{definition}
\begin{proof}
  This is the declared model definition of signed Charge In Coulombs.
\end{proof}
\begin{definition}[coulomb Constant In Joule Meters Per Coulomb Squared]
  \label{def:physics:phyx-mini-0636:phyxminiproblems-problemphyxmini0636-coulombconstantinjoulemeterspercoulombsquared}
  \lean{PhyXMiniProblems.ProblemPhyXMini0636.coulombConstantInJouleMetersPerCoulombSquared}
  \uses{def:physics:phyx-mini-0636:phyxminiproblems-problemphyxmini0636-coulombconstantquantity}
  Read a dimensionful value of Coulomb's constant in J m / C²
\end{definition}
\begin{proof}
  This is the declared model definition of coulomb Constant In Joule Meters Per Coulomb Squared.
\end{proof}
\begin{definition}[energy In Joules]
  \label{def:physics:phyx-mini-0636:phyxminiproblems-problemphyxmini0636-energyinjoules}
  \lean{PhyXMiniProblems.ProblemPhyXMini0636.energyInJoules}
  Read a signed physical energy in coherent-SI joules
\end{definition}
\begin{proof}
  This is the declared model definition of energy In Joules.
\end{proof}
\begin{definition}[elementary Charge In Coulombs]
  \label{def:physics:phyx-mini-0636:phyxminiproblems-problemphyxmini0636-elementarychargeincoulombs}
  \lean{PhyXMiniProblems.ProblemPhyXMini0636.elementaryChargeInCoulombs}
  Physlib's positive elementary-charge unit expressed in coulombs
\end{definition}
\begin{proof}
  This is the declared model definition of elementary Charge In Coulombs.
\end{proof}
\begin{definition}[signed Charge In Elementary Charges]
  \label{def:physics:phyx-mini-0636:phyxminiproblems-problemphyxmini0636-signedchargeinelementarycharges}
  \lean{PhyXMiniProblems.ProblemPhyXMini0636.signedChargeInElementaryCharges}
  \uses{def:physics:phyx-mini-0636:phyxminiproblems-problemphyxmini0636-signedchargequantity, def:physics:phyx-mini-0636:phyxminiproblems-problemphyxmini0636-signedchargeincoulombs, def:physics:phyx-mini-0636:phyxminiproblems-problemphyxmini0636-elementarychargeincoulombs}
  Read a signed charge in units of the positive elementary charge
\end{definition}
\begin{proof}
  This is the declared model definition of signed Charge In Elementary Charges.
\end{proof}
\begin{definition}[electron Volt In Joules]
  \label{def:physics:phyx-mini-0636:phyxminiproblems-problemphyxmini0636-electronvoltinjoules}
  \lean{PhyXMiniProblems.ProblemPhyXMini0636.electronVoltInJoules}
  \uses{def:physics:phyx-mini-0636:phyxminiproblems-problemphyxmini0636-energyinjoules}
  The joule readout of Physlib's dimensionful electron volt
\end{definition}
\begin{proof}
  This is the declared model definition of electron Volt In Joules.
\end{proof}
\begin{definition}[kilo Electron Volt In Joules]
  \label{def:physics:phyx-mini-0636:phyxminiproblems-problemphyxmini0636-kiloelectronvoltinjoules}
  \lean{PhyXMiniProblems.ProblemPhyXMini0636.kiloElectronVoltInJoules}
  \uses{def:physics:phyx-mini-0636:phyxminiproblems-problemphyxmini0636-electronvoltinjoules}
  One kiloelectronvolt expressed in joules
\end{definition}
\begin{proof}
  This is the declared model definition of kilo Electron Volt In Joules.
\end{proof}
\begin{definition}[mega Electron Volt In Joules]
  \label{def:physics:phyx-mini-0636:phyxminiproblems-problemphyxmini0636-megaelectronvoltinjoules}
  \lean{PhyXMiniProblems.ProblemPhyXMini0636.megaElectronVoltInJoules}
  \uses{def:physics:phyx-mini-0636:phyxminiproblems-problemphyxmini0636-electronvoltinjoules}
  One megaelectronvolt expressed in joules
\end{definition}
\begin{proof}
  This is the declared model definition of mega Electron Volt In Joules.
\end{proof}
\begin{definition}[energy In Kilo Electron Volts]
  \label{def:physics:phyx-mini-0636:phyxminiproblems-problemphyxmini0636-energyinkiloelectronvolts}
  \lean{PhyXMiniProblems.ProblemPhyXMini0636.energyInKiloElectronVolts}
  \uses{def:physics:phyx-mini-0636:phyxminiproblems-problemphyxmini0636-energyinjoules, def:physics:phyx-mini-0636:phyxminiproblems-problemphyxmini0636-kiloelectronvoltinjoules}
  Read a physical energy in kiloelectronvolts
\end{definition}
\begin{proof}
  This is the declared model definition of energy In Kilo Electron Volts.
\end{proof}
\begin{definition}[energy In Mega Electron Volts]
  \label{def:physics:phyx-mini-0636:phyxminiproblems-problemphyxmini0636-energyinmegaelectronvolts}
  \lean{PhyXMiniProblems.ProblemPhyXMini0636.energyInMegaElectronVolts}
  \uses{def:physics:phyx-mini-0636:phyxminiproblems-problemphyxmini0636-energyinjoules, def:physics:phyx-mini-0636:phyxminiproblems-problemphyxmini0636-megaelectronvoltinjoules}
  Read a physical energy in megaelectronvolts
\end{definition}
\begin{proof}
  This is the declared model definition of energy In Mega Electron Volts.
\end{proof}
\begin{definition}[Interaction State]
  \label{def:physics:phyx-mini-0636:phyxminiproblems-problemphyxmini0636-interactionstate}
  \lean{PhyXMiniProblems.ProblemPhyXMini0636.InteractionState}
  The two states labelled Before and After in the supplied figure
\end{definition}
\begin{proof}
  This is the declared model definition of Interaction State.
\end{proof}
\begin{definition}[Radial Region]
  \label{def:physics:phyx-mini-0636:phyxminiproblems-problemphyxmini0636-radialregion}
  \lean{PhyXMiniProblems.ProblemPhyXMini0636.RadialRegion}
  Nuclear regions occupied by the emitted beta particle in the model
\end{definition}
\begin{proof}
  This is the declared model definition of Radial Region.
\end{proof}
\begin{definition}[Decay Mode]
  \label{def:physics:phyx-mini-0636:phyxminiproblems-problemphyxmini0636-decaymode}
  \lean{PhyXMiniProblems.ProblemPhyXMini0636.DecayMode}
  The radioactive mode stated in the problem
\end{definition}
\begin{proof}
  This is the declared model definition of Decay Mode.
\end{proof}
\begin{definition}[Particle Species]
  \label{def:physics:phyx-mini-0636:phyxminiproblems-problemphyxmini0636-particlespecies}
  \lean{PhyXMiniProblems.ProblemPhyXMini0636.ParticleSpecies}
  Particle species distinguished by this model
\end{definition}
\begin{proof}
  This is the declared model definition of Particle Species.
\end{proof}
\begin{definition}[Horizontal Direction]
  \label{def:physics:phyx-mini-0636:phyxminiproblems-problemphyxmini0636-horizontaldirection}
  \lean{PhyXMiniProblems.ProblemPhyXMini0636.HorizontalDirection}
  Qualitative direction of the arrows drawn beside the beta particle
\end{definition}
\begin{proof}
  This is the declared model definition of Horizontal Direction.
\end{proof}
\begin{definition}[Cesium Beta Decay Figure]
  \label{def:physics:phyx-mini-0636:phyxminiproblems-problemphyxmini0636-cesiumbetadecayfigure}
  \lean{PhyXMiniProblems.ProblemPhyXMini0636.CesiumBetaDecayFigure}
  \uses{def:physics:phyx-mini-0636:phyxminiproblems-problemphyxmini0636-interactionstate, def:physics:phyx-mini-0636:phyxminiproblems-problemphyxmini0636-radialregion, def:physics:phyx-mini-0636:phyxminiproblems-problemphyxmini0636-horizontaldirection}
  Literal presentation data from image 636. The numerical fields are printed readouts; their calibration to physical quantities is stated separately
\end{definition}
\begin{proof}
  This is the declared model definition of Cesium Beta Decay Figure.
\end{proof}
\begin{definition}[Cesium Beta Ejection Setup]
  \label{def:physics:phyx-mini-0636:phyxminiproblems-problemphyxmini0636-cesiumbetaejectionsetup}
  \lean{PhyXMiniProblems.ProblemPhyXMini0636.CesiumBetaEjectionSetup}
  \uses{def:physics:phyx-mini-0636:phyxminiproblems-problemphyxmini0636-lengthquantity, def:physics:phyx-mini-0636:phyxminiproblems-problemphyxmini0636-signedchargequantity, def:physics:phyx-mini-0636:phyxminiproblems-problemphyxmini0636-coulombconstantquantity, def:physics:phyx-mini-0636:phyxminiproblems-problemphyxmini0636-interactionstate, def:physics:phyx-mini-0636:phyxminiproblems-problemphyxmini0636-radialregion, def:physics:phyx-mini-0636:phyxminiproblems-problemphyxmini0636-decaymode, def:physics:phyx-mini-0636:phyxminiproblems-problemphyxmini0636-particlespecies, def:physics:phyx-mini-0636:phyxminiproblems-problemphyxmini0636-cesiumbetadecayfigure}
  Independent physical data for the beta-particle/nucleus interaction. The before-state kinetic energy is stored as an observable and is constrained only through the general laws below; it is not defined from 13.1 MeV
\end{definition}
\begin{proof}
  This is the declared model definition of Cesium Beta Ejection Setup.
\end{proof}
\begin{definition}[Matches Cesium137 Beta Decay Scenario]
  \label{def:physics:phyx-mini-0636:phyxminiproblems-problemphyxmini0636-matchescesium137betadecayscenario}
  \lean{PhyXMiniProblems.ProblemPhyXMini0636.MatchesCesium137BetaDecayScenario}
  \uses{def:physics:phyx-mini-0636:phyxminiproblems-problemphyxmini0636-lengthinfemtometers, def:physics:phyx-mini-0636:phyxminiproblems-problemphyxmini0636-cesiumbetaejectionsetup}
  Cesium-137, atomic number 55, undergoing beta-minus decay
\end{definition}
\begin{proof}
  This is the declared model definition of Matches Cesium137 Beta Decay Scenario.
\end{proof}
\begin{definition}[Matches Supplied Cesium Beta Figure]
  \label{def:physics:phyx-mini-0636:phyxminiproblems-problemphyxmini0636-matchessuppliedcesiumbetafigure}
  \lean{PhyXMiniProblems.ProblemPhyXMini0636.MatchesSuppliedCesiumBetaFigure}
  \uses{def:physics:phyx-mini-0636:phyxminiproblems-problemphyxmini0636-cesiumbetaejectionsetup}
  All unambiguous symbols, words, arrows, and numbers visible in image 636
\end{definition}
\begin{proof}
  This is the declared model definition of Matches Supplied Cesium Beta Figure.
\end{proof}
\begin{definition}[Uses Supplied Figure Calibration]
  \label{def:physics:phyx-mini-0636:phyxminiproblems-problemphyxmini0636-usessuppliedfigurecalibration}
  \lean{PhyXMiniProblems.ProblemPhyXMini0636.UsesSuppliedFigureCalibration}
  \uses{def:physics:phyx-mini-0636:phyxminiproblems-problemphyxmini0636-lengthinfemtometers, def:physics:phyx-mini-0636:phyxminiproblems-problemphyxmini0636-signedchargeinelementarycharges, def:physics:phyx-mini-0636:phyxminiproblems-problemphyxmini0636-energyinkiloelectronvolts, def:physics:phyx-mini-0636:phyxminiproblems-problemphyxmini0636-cesiumbetaejectionsetup}
  Calibration of the literal figure readouts to the independent physical quantities. Only the final observed energy and the two charges are supplied; the requested before-state kinetic energy is absent
\end{definition}
\begin{proof}
  This is the declared model definition of Uses Supplied Figure Calibration.
\end{proof}
\begin{definition}[Has Nuclear Surface Geometry]
  \label{def:physics:phyx-mini-0636:phyxminiproblems-problemphyxmini0636-hasnuclearsurfacegeometry}
  \lean{PhyXMiniProblems.ProblemPhyXMini0636.HasNuclearSurfaceGeometry}
  \uses{def:physics:phyx-mini-0636:phyxminiproblems-problemphyxmini0636-lengthinmeters, def:physics:phyx-mini-0636:phyxminiproblems-problemphyxmini0636-cesiumbetaejectionsetup}
  The diameter, radius, and surface separation describe the same nucleus
\end{definition}
\begin{proof}
  This is the declared model definition of Has Nuclear Surface Geometry.
\end{proof}
\begin{definition}[Has Physical Cesium Beta Parameters]
  \label{def:physics:phyx-mini-0636:phyxminiproblems-problemphyxmini0636-hasphysicalcesiumbetaparameters}
  \lean{PhyXMiniProblems.ProblemPhyXMini0636.HasPhysicalCesiumBetaParameters}
  \uses{def:physics:phyx-mini-0636:phyxminiproblems-problemphyxmini0636-lengthinmeters, def:physics:phyx-mini-0636:phyxminiproblems-problemphyxmini0636-coulombconstantinjoulemeterspercoulombsquared, def:physics:phyx-mini-0636:phyxminiproblems-problemphyxmini0636-energyinjoules, def:physics:phyx-mini-0636:phyxminiproblems-problemphyxmini0636-elementarychargeincoulombs, def:physics:phyx-mini-0636:phyxminiproblems-problemphyxmini0636-electronvoltinjoules, def:physics:phyx-mini-0636:phyxminiproblems-problemphyxmini0636-cesiumbetaejectionsetup}
  Positivity and sign conditions for the physical parameters
\end{definition}
\begin{proof}
  This is the declared model definition of Has Physical Cesium Beta Parameters.
\end{proof}
\begin{definition}[Uses Vacuum Coulomb Constant]
  \label{def:physics:phyx-mini-0636:phyxminiproblems-problemphyxmini0636-usesvacuumcoulombconstant}
  \lean{PhyXMiniProblems.ProblemPhyXMini0636.UsesVacuumCoulombConstant}
  \uses{def:physics:phyx-mini-0636:phyxminiproblems-problemphyxmini0636-coulombconstantinjoulemeterspercoulombsquared, def:physics:phyx-mini-0636:phyxminiproblems-problemphyxmini0636-cesiumbetaejectionsetup}
  The dimensionful Coulomb constant is calibrated to Physlib's electromagnetic system and to the standard vacuum SI readout. This supplies a physical constant, not the requested kinetic-energy answer
\end{definition}
\begin{proof}
  This is the declared model definition of Uses Vacuum Coulomb Constant.
\end{proof}
\begin{definition}[Satisfies Ideal Electrostatic Ejection Laws]
  \label{def:physics:phyx-mini-0636:phyxminiproblems-problemphyxmini0636-satisfiesidealelectrostaticejectionlaws}
  \lean{PhyXMiniProblems.ProblemPhyXMini0636.SatisfiesIdealElectrostaticEjectionLaws}
  \uses{def:physics:phyx-mini-0636:phyxminiproblems-problemphyxmini0636-lengthinmeters, def:physics:phyx-mini-0636:phyxminiproblems-problemphyxmini0636-signedchargeincoulombs, def:physics:phyx-mini-0636:phyxminiproblems-problemphyxmini0636-coulombconstantinjoulemeterspercoulombsquared, def:physics:phyx-mini-0636:phyxminiproblems-problemphyxmini0636-energyinjoules, def:physics:phyx-mini-0636:phyxminiproblems-problemphyxmini0636-cesiumbetaejectionsetup}
  The beta particle starts one nuclear radius from the centre, so its initial electrostatic potential energy is the attractive point-charge Coulomb energy. Mechanical energy K + U is conserved between the two displayed states. Neither relation contains a numerical value for the initial kinetic energy
\end{definition}
\begin{proof}
  This is the declared model definition of Satisfies Ideal Electrostatic Ejection Laws.
\end{proof}
\begin{definition}[Answer Choice]
  \label{def:physics:phyx-mini-0636:phyxminiproblems-problemphyxmini0636-answerchoice}
  \lean{PhyXMiniProblems.ProblemPhyXMini0636.AnswerChoice}
  Labels of the four answers printed with the problem
\end{definition}
\begin{proof}
  This is the declared model definition of Answer Choice.
\end{proof}
\begin{definition}[ejection Kinetic Energy In Mega Electron Volts]
  \label{def:physics:phyx-mini-0636:phyxminiproblems-problemphyxmini0636-answerchoice-ejectionkineticenergyinmegaelectronvolts}
  \lean{PhyXMiniProblems.ProblemPhyXMini0636.AnswerChoice.ejectionKineticEnergyInMegaElectronVolts}
  \uses{def:physics:phyx-mini-0636:phyxminiproblems-problemphyxmini0636-answerchoice}
  Ejection kinetic energy printed beside each answer, in megaelectronvolts
\end{definition}
\begin{proof}
  This is the declared model definition of ejection Kinetic Energy In Mega Electron Volts.
\end{proof}
\begin{definition}[displayed Energy Rounding Tolerance In Mega Electron Volts]
  \label{def:physics:phyx-mini-0636:phyxminiproblems-problemphyxmini0636-displayedenergyroundingtoleranceinmegaelectronvolts}
  \lean{PhyXMiniProblems.ProblemPhyXMini0636.displayedEnergyRoundingToleranceInMegaElectronVolts}
  Half of the 0.1 MeV precision displayed by the answer values
\end{definition}
\begin{proof}
  This is the declared model definition of displayed Energy Rounding Tolerance In Mega Electron Volts.
\end{proof}
\begin{definition}[Matches Displayed Ejection Energy]
  \label{def:physics:phyx-mini-0636:phyxminiproblems-problemphyxmini0636-matchesdisplayedejectionenergy}
  \lean{PhyXMiniProblems.ProblemPhyXMini0636.MatchesDisplayedEjectionEnergy}
  \uses{def:physics:phyx-mini-0636:phyxminiproblems-problemphyxmini0636-energyinmegaelectronvolts, def:physics:phyx-mini-0636:phyxminiproblems-problemphyxmini0636-cesiumbetaejectionsetup, def:physics:phyx-mini-0636:phyxminiproblems-problemphyxmini0636-answerchoice, def:physics:phyx-mini-0636:phyxminiproblems-problemphyxmini0636-displayedenergyroundingtoleranceinmegaelectronvolts}
  The physical ejection energy rounds to the value printed for a choice
\end{definition}
\begin{proof}
  This is the declared model definition of Matches Displayed Ejection Energy.
\end{proof}
\begin{definition}[Is Unique Matching Displayed Ejection Energy]
  \label{def:physics:phyx-mini-0636:phyxminiproblems-problemphyxmini0636-isuniquematchingdisplayedejectionenergy}
  \lean{PhyXMiniProblems.ProblemPhyXMini0636.IsUniqueMatchingDisplayedEjectionEnergy}
  \uses{def:physics:phyx-mini-0636:phyxminiproblems-problemphyxmini0636-cesiumbetaejectionsetup, def:physics:phyx-mini-0636:phyxminiproblems-problemphyxmini0636-answerchoice, def:physics:phyx-mini-0636:phyxminiproblems-problemphyxmini0636-matchesdisplayedejectionenergy}
  Exactly one printed choice matches at the displayed precision
\end{definition}
\begin{proof}
  This is the declared model definition of Is Unique Matching Displayed Ejection Energy.
\end{proof}
\begin{definition}[recorded Dataset Answer]
  \label{def:physics:phyx-mini-0636:phyxminiproblems-problemphyxmini0636-recordeddatasetanswer}
  \lean{PhyXMiniProblems.ProblemPhyXMini0636.recordedDatasetAnswer}
  \uses{def:physics:phyx-mini-0636:phyxminiproblems-problemphyxmini0636-answerchoice}
  The source dataset's recorded answer, retained only as metadata
\end{definition}
\begin{proof}
  This is the declared model definition of recorded Dataset Answer.
\end{proof}
\begin{lemma}[ejection Kinetic Energy coulomb Balance]
  \label{lem:physics:phyx-mini-0636:phyxminiproblems-problemphyxmini0636-ejectionkineticenergy-coulombbalance}
  \lean{PhyXMiniProblems.ProblemPhyXMini0636.ejectionKineticEnergy_coulombBalance}
  \uses{def:physics:phyx-mini-0636:phyxminiproblems-problemphyxmini0636-lengthinmeters, def:physics:phyx-mini-0636:phyxminiproblems-problemphyxmini0636-signedchargeincoulombs, def:physics:phyx-mini-0636:phyxminiproblems-problemphyxmini0636-coulombconstantinjoulemeterspercoulombsquared, def:physics:phyx-mini-0636:phyxminiproblems-problemphyxmini0636-energyinjoules, def:physics:phyx-mini-0636:phyxminiproblems-problemphyxmini0636-cesiumbetaejectionsetup, def:physics:phyx-mini-0636:phyxminiproblems-problemphyxmini0636-matchessuppliedcesiumbetafigure, def:physics:phyx-mini-0636:phyxminiproblems-problemphyxmini0636-usessuppliedfigurecalibration, def:physics:phyx-mini-0636:phyxminiproblems-problemphyxmini0636-satisfiesidealelectrostaticejectionlaws}
  Conservation of energy and the surface Coulomb law first give the exact coherent-SI balance, before any numerical constants are substituted
\end{lemma}
\begin{proof}
  The conclusion follows from the governing relations and figure readouts named in the statement of ejection Kinetic Energy coulomb Balance.
\end{proof}
\begin{lemma}[ejection Kinetic Energy eq coulomb Corrected Final Energy]
  \label{lem:physics:phyx-mini-0636:phyxminiproblems-problemphyxmini0636-ejectionkineticenergy-eq-coulombcorrectedfinalenergy}
  \lean{PhyXMiniProblems.ProblemPhyXMini0636.ejectionKineticEnergy_eq_coulombCorrectedFinalEnergy}
  \uses{def:physics:phyx-mini-0636:phyxminiproblems-problemphyxmini0636-lengthinmeters, def:physics:phyx-mini-0636:phyxminiproblems-problemphyxmini0636-coulombconstantinjoulemeterspercoulombsquared, def:physics:phyx-mini-0636:phyxminiproblems-problemphyxmini0636-elementarychargeincoulombs, def:physics:phyx-mini-0636:phyxminiproblems-problemphyxmini0636-megaelectronvoltinjoules, def:physics:phyx-mini-0636:phyxminiproblems-problemphyxmini0636-energyinmegaelectronvolts, def:physics:phyx-mini-0636:phyxminiproblems-problemphyxmini0636-cesiumbetaejectionsetup, def:physics:phyx-mini-0636:phyxminiproblems-problemphyxmini0636-matchessuppliedcesiumbetafigure, def:physics:phyx-mini-0636:phyxminiproblems-problemphyxmini0636-usessuppliedfigurecalibration, def:physics:phyx-mini-0636:phyxminiproblems-problemphyxmini0636-hasnuclearsurfacegeometry, def:physics:phyx-mini-0636:phyxminiproblems-problemphyxmini0636-hasphysicalcesiumbetaparameters, def:physics:phyx-mini-0636:phyxminiproblems-problemphyxmini0636-satisfiesidealelectrostaticejectionlaws}
  After inserting the figure charges, radius, and final 300 keV, the ejection energy is 0.3 MeV plus the positive magnitude of the attractive surface Coulomb potential
\end{lemma}
\begin{proof}
  The conclusion follows from the governing relations and figure readouts named in the statement of ejection Kinetic Energy eq coulomb Corrected Final Energy.
\end{proof}
% --- Archon declaration topology end ---
% --- Archon physics formalization source end ---
